\chapter{What Is the World Made Of?}
\section{A Bicycle with Every Part Replaced}
First, pick up an object: an old bicycle.

Not a new bike, but one used for more than a decade. The rubber grips are polished with wear, the saddle springs squeak when you sit, the chain has been replaced three times, the pedals twice, and a wheel rim once after a collision bent it. Even the frame has a welded repair. It originally belonged to your grandfather, then your uncle; now you ride it.

Here is a question that sounds rather dull: is this still the original bicycle?

You will probably say of course; it has always been this one. But suppose someone assembles the discarded chain, pedals, and wheel rim from the garage into another bike, puts it before you, and says, ``This is the original. These are its original parts.'' How would you answer?

Do not rush. This question is much larger than it seems. It concerns not only bicycles but everything. Is a river the same river after all its water changes? Is a class the same class after successive generations of students? As your cells continually renew, are you still the person you were seven years ago? A stone, a thought, a shower, a friendship: are these one kind of existence or fundamentally different kinds?

Pile these questions together and you have this chapter's problem: \textbf{what is the world actually made of?} It is among humanity's oldest questions, seriously answered in almost every civilization, with entirely different answers. The philosophical branch studying it is \textbf{metaphysics}. Do not fear the name. Its basic concern is simple: physics studies how things in the world move; metaphysics asks an earlier question: what kinds of things exist at all? And what does thing itself mean?

\section{An Afternoon in a Bicycle Repair Shop}
Now enter a scene with me.

Time: a summer afternoon, after four, with the sun still fierce. Place: the street behind your apartment complex, at a repair shop operating for more than twenty years. It is tiny. A dozen old inner tubes hang outside like black sausages. Oil and rubber mingle with the burnt-metal smell of cutting from the hardware shop next door. A radio quietly plays traditional storytelling.

The repairer, Mr. Geng, is in his fifties, with oily grime permanently embedded in his hands. He crouches to replace an old bicycle's bottom bracket, the rotating axle between its pedals. Its gray-haired owner watches from a small folding stool in the shade.

``Mr. Geng,'' the old man says, ``this bike has been with me thirty years. Apart from the frame, I think everything's been replaced.''

Without looking up, Geng replies, ``The frame won't last much longer either. The front fork's been welded and the rear has cracked.''

``Then tell me,'' the owner smiles, ``is it still the original bike?''

Geng stops, thinks seriously, and says something that makes a waiting secondary-school student look up:

``Replace every part and it's still itself. I handle it every day. I recognize it.''

The owner persists: ``What exactly do you recognize? None of the parts are original.''

``It's...this bike.'' Geng pats the frame as though that answers everything.

There is no professor or textbook this afternoon, only a craftsman and an old man, yet they debate one of Western philosophy's most famous problems. Geng recognizes something \textbf{not located in the parts}. The owner asks what entitles a bike with nothing original remaining to be itself. One recognizes an intangible this one, the other tangible components.

Neither persuades the other because their understandings of what a bicycle is differ fundamentally. We will now enlarge this little dispute to the whole world.

\section{The Question: What Kinds of Things Exist?}
Before answering, see the conflict clearly.

Look around and list existing things: table, phone, air, sun. These seem straightforward, occupying space and tangible. Philosophers call them \textbf{objects} or \textbf{matter}, things composed of material and occupying space.

But the list soon encounters trouble.

Does my headache count as a thing? It occupies no space, and others cannot see or touch it, yet it is real enough to make me wince. This is \textbf{mental} existence: sensations, thoughts, dreams, pain.

Does yesterday's match count? It is no object you can put in a box. It has a beginning, course, and ending. This is an \textbf{event}, something that happens.

River water flowing, an economy growing, you growing up: these resemble \textbf{processes}, changes unfolding rather than single completed happenings.

I am taller than you, the school stands beside the post office, this caused that: these are \textbf{relations}. Neither objects nor sensations, they nevertheless prevent the world becoming disconnected fragments. Remove taller, beside, and because, and what remains?

We now have at least five candidates: matter, mind, events, processes, relations. Philosophers also use three old conceptual pairs to inspect the list. Meet them now; they will recur.

\begin{itemize}
\item \textbf{Substance and property}: a substance exists independently, like this bicycle. A property, its redness or heaviness, depends on a substance. Redness cannot walk down the street alone; it must belong to something.
\item \textbf{Part and whole}: wheels, frame, and chain are parts; the bicycle is the whole. Does a whole simply equal its parts added together? Are scattered components already a bike? If not, what additional being-a-bike exists, and where beyond the parts?
\item \textbf{Identity and change}: things continually change, yet we call them the same. What justifies that same?
\end{itemize}
Here is the real problem:

\textbf{Are all five kinds equally real, or is one fundamental and the rest its variations?}

One side says only matter is real: headaches are neurons firing, matches are bodies moving, relations are ways of speaking. Ultimately there is matter and what happens to it. This is \textbf{materialism}.

Another says matter and mind differ fundamentally and neither reduces to the other. The brain is material, but pain's feeling is not a lump of matter; it exists on another level. This is \textbf{dualism}.

A third reverses matters: mind is fundamental. You never encounter matter outside all perception. Every table you know is seen or touched. Perception comes first, then what you call the world. This is \textbf{idealism}.

For over two millennia these sides have disputed not only the world's constituents but whether an object survives replacement of every part. Identity depends on what you think it \textbf{fundamentally is}: material, shape, continuous history, or simply a name.

This is the problem of \textbf{identity}, forever at odds with change. Things change, yet remain the same in our speech. Why?

\section{Three Positions: First Hear Them Out}
The honest response to competing positions is not choosing a favorite first, but completing every side's reasoning, even if you ultimately disagree. Let us practice making another side's full case.

\textbf{Position one: materialism, matter and its motion alone.}

This idea is extraordinarily old. More than 2,400 years ago, Democritus held that indivisible particles, atoms, combine and move in a void to form everything. Smells are atoms entering noses; thoughts are atoms colliding in heads. A joke? It is an ancestor of modern science. Today's materialists, more often called physicalists, offer strong reasons.

First, science's history repeatedly explains supposedly mysterious things. Vital force became cellular chemistry; supernatural disease became bacteria and viruses. Claims that matter could never explain something repeatedly failed. Why should mind be the exception?

Second, physical causation is closed. Raising a hand follows a complete chain of neural firing and muscular contraction. If thoughts are independently nonphysical, how can they push neurons? There seems no place to intervene.

Third, materialism is simplest. One kind of thing explains the world without inventing invisible, unmeasurable, obscure mental stuff.

\textbf{Position two: dualism, fundamentally different matter and mind.}

Hear this case fully too. Its best-known advocate, seventeenth-century French philosopher Descartes, begins with reasoning anyone can attempt. I may doubt tables, my body, or an external world that could be a dream. But I cannot doubt the very occurrence of doubting or thinking. A doubting thought necessarily exists. Hence his Discourse on Method states:

\begin{quote}
I think, therefore I am. (Je pense, donc je suis)
\end{quote}
He concludes that thinking's existence is more certain than anything material. The body might not exist, perhaps being dreamed, but I as thinker cannot be nonexistent. Thus mind and matter differ: matter occupies space and divides, while mind occupies none and a thought cannot be cut in half.

Dualism retains force today. You know what pain feels like, while even the best neuroscientist sees firing patterns on a brain scan, not pain itself. Feeling has an internally experienced aspect inaccessible to external material observation. Philosophers call this \textbf{subjective experience}: the what-it-feels-like in being you.

\textbf{Position three: idealism, existence as being perceived.}

This position is easiest to ridicule, so particularly needs its full argument. In A Treatise Concerning the Principles of Human Knowledge, eighteenth-century Irish philosopher Berkeley argues, broadly, that being is being perceived. His reasoning is tricky: imagine a table perceived by nobody. Even your imagined table is imagined by you. You cannot separate the table from awareness and grasp an entirely unperceived table. What warrants asserting a world wholly independent of all perception?

Ming thinker Wang Yangming offers an almost parallel account. Instructions for Practical Living records someone pointing to a mountain flower and asking how its unobserved blooming and falling relate to the mind if nothing exists outside mind. Wang replies:

\begin{quote}
Before you look at this flower, it and your mind rest together in stillness. When you look, its colors become clear at once. Thus you know it is not outside your mind.
\end{quote}
These positions need not mean tables do not exist or cannot hurt when kicked, a common misreading. Samuel Johnson famously responded to Berkeley by kicking a large stone and announcing his refutation. Berkeley could agree entirely that it hurt. The dispute concerns not hardness but whether hardness has any footing apart from perception.

Each camp also argues internally. Radical materialists call mind illusory; moderates accept feelings while expecting scientific explanation. Song thinker Zhang Zai followed another route: in Correcting Ignorance, the great void is qi, with things forming through qi's gathering and dispersal. Neither Western atomism nor idealism, this treats objects as concentrations of a continuous medium. Indian Buddhist Yogacara likewise links all experience to consciousness's activity. \textbf{No tradition speaks with only one voice}. Remembering that matters more than school names.

\section[Inventory, Then Test Identity]{Thinking Bridge: Inventory, Then Test Identity}
Can a portable tool replace intuition in this immense question? Yes, in two simple but powerful steps.

\textbf{Step one: list what exists and classify each item.}

In any philosophical or everyday dispute, list its things and ask whether each is an object, dependent property such as redness or hardness, event, process, relation, or mental state.

Disputes often treat unlike categories alike. Does an old school still exist? One person means its demolished building; another means relationships and traditions persisting after demolition. Three days of argument cannot settle different categories of school. Classification alone removes much conflict.

\textbf{Step two: ask each item's criterion of identity.}

The criterion is the ruler deciding when something stays the same or becomes another. \textbf{Different kinds need different rulers.}

\begin{itemize}
\item A sand pile's identity rests chiefly on material: replace half and it is practically another pile.
\item A football team's does not. Players, stadium, and badge can all change while inheritance preserves the same team. Its ruler is \textbf{continuous organization and history}.
\item A river's ruler is its enduring channel and drainage structure. Constantly replaced water does not make a new Yellow River each second.
\item A musical piece has a looser ruler. Piano or suona, faster or slower, retained melodic structure preserves it, almost independently of material.
\end{itemize}
The bicycle puzzle is difficult not through obscurity but because it straddles rulers: material collection and functional, historical artifact. One side measures sand, the other a team. Agreement cannot follow incompatible measures.

Use this today when people dispute whether an updated game or relocated school remains itself. Before choosing sides, ask: \textbf{what kind of thing is it, and which ruler measures its identity?}

\section{Reading Evidence: Theseus's Ship}
The opening bicycle has a two-thousand-year-old predecessor.

Around the turn of the first and second centuries CE, Plutarch's Life of Theseus in Parallel Lives records, in paraphrase, Athens preserving the hero's returning ship. Decayed planks were replaced with sound ones, one, two, ten, until centuries later none remained original. It became a living philosophical dispute: one side held the ship identical, another denied it.

Centuries later, English philosopher Hobbes added a devastating twist. Collect every discarded plank and reassemble them. Two ships now exist: continuous repairs with new wood, and later reassembly with original wood. \textbf{Which is Theseus's ship?}

His contemporary Locke addresses such questions in An Essay Concerning Human Understanding, here paraphrased. Identity standards differ by kind: material for a stone, continuous living organization for an oak growing from seed despite repeated replacement, and continuous memory and consciousness for a person. Remembered deeds count as yours. Three thinkers and three rulers anticipate our tool across two millennia.

Look eastward too. Not far from Plutarch's era, Heraclitus supplied a sharper image. The familiar inability to enter one river twice is a later paraphrase, not his exact words: the second entry meets different water; all things similarly flow. In Zhuangzi's Warring States ``Autumn Floods,'' two thinkers watch fish from a bridge:

\begin{quote}
Zhuangzi said, ``The minnows swim at ease: this is fish's joy.'' Huizi replied, ``You are not a fish. How do you know its joy?'' Zhuangzi answered, ``You are not me. How do you know I do not know its joy?''
\end{quote}
Beneath apparent banter lie our two problems: how can Zhuangzi assert a fish's mental joy, and what permits constantly changing rivers and fish to remain that river and group?

Across millennia and civilizations, ships, rivers, fish, trees, and people repeatedly \textbf{test our most basic classifications of things}. Metaphysics is not airy speculation, but unpacking the objects, properties, events, processes, relations, and minds we use daily without inspection.

\section{One Ship, Three Answers}
With inventory and sources, compare answers. Distinguish replaced planks as evidence; differing conceptions of things as explanation; Athenians preserving what they regarded as Theseus's ship as participants' understanding; and the better answer as your evaluation. Compare explanations.

\textbf{Answer one: material. The original planks make the original ship.}

Reason: a thing consists of its material; lose that and lose the thing. The new-wood ship is a successor or copy. Museums tacitly honor this intuition: visitors pay to see the original Mona Lisa, not an indistinguishable reproduction, revealing a premium on original material.

The limit is fatal: you would no longer be yourself seven years earlier; the oak would not remain its acorn's tree. Even melting ice remains the same water in our speech, using continuous change rather than material as our ruler. Material alone fails everyday usage, working best for particular authenticity questions.

\textbf{Answer two: continuity. The continuously repaired ship is original.}

Reason: identity rests not on a batch of material but an unbroken history of shape, structure, function, and use. Each new plank enters existing organization like a new student entering an old class. Identity passes like a torch, explaining oaks, rivers, teams, and you.

Its limit is Hobbes's rival ship. If dismantling and storage interrupted its history, does all that original wood lose qualification? Many cannot accept the original's whole material no longer being the ship. Continuity's boundary is also vague: where does conservation become reconstruction? Would a wholly rebuilt Old Summer Palace still be that palace?

\textbf{Answer three: convention. Identity has no hidden worldly answer, only agreed usage.}

Reason: identity is not a factory label but a convenient linguistic boundary. Matter, events, and processes flow; calling something the same ship helps discussion, as do nation, class, and market price. The puzzle asks how we intend to use same ship, not which answer is objectively correct.

Its explanatory strength borders on evasion. If all identity is conventional, your concern about being the one who wakes tomorrow becomes dictionary usage. That worry seems deeper than changed definitions. Translating every existential question into language can annul rather than solve it.

Each answer holds something: original material's weight, persistence through change, human classification. Each misses a corner. This is not bland compromise, but a method: \textbf{when an answer seems unassailable, look for what it discarded first}.

\section[What If Things Are Processes?]{Deeper Challenge: What If Things Are Processes?}
Our strongest theory challenges an unnoticed assumption: \textbf{the world consists of things}.

So far, the world resembles a room containing ships, trees, and people to which changes happen. Objects star; change is their performance.

In Process and Reality, early-twentieth-century philosopher Whitehead reverses this. \textbf{Events and processes} are fundamental; objects are relatively stable, slowly changing processes mistaken for fixed things. Flame is the best example: fuel enters, gases leave, no moment's flame material belongs to the next, yet the flame exists. Through this lens, \textbf{everything is flame}, burning at different speeds. A stone's metaphorical burning spans hundreds of millions of years, too slowly for us to notice.

Around then, physics approached a related picture. Special relativity weaves time and space into four-dimensional spacetime. Objects become not three-dimensional blocks moving through time but four-dimensional worms extended along it. Your present self is a \textbf{temporal slice}, like a film frame; yesterday, today, and tomorrow are cross-sections. Ship identity becomes whether slices have sufficiently close continuity, not whether changed material qualifies as original. Philosophers call complete presence at each moment endurantism, and possession of temporal parts perdurantism. Their dispute continues.

Process thinking has an ancient Eastern relative. Buddhism observes conditioned things arising and passing moment by moment, without fixed self-nature. The Diamond Sutra concludes:

\begin{quote}
All conditioned things are like dreams, illusions, bubbles, shadows, like dew and lightning: see them thus.
\end{quote}
The point is not nonexistence but temporary rather than permanent residence. Lightning exists and lights the sky, yet no fixed lightning object can be found, only a moment's discharge.

Keep this chapter's repeated balance. Buddhist thought is introduced to understand a worldview influencing billions, not require belief. Whitehead and spacetime offer another way to see, not abolish ordinary speech about your bike or school. At boundary cases, however, deeper pictures may help when ordinary descriptions fail.

Mark the theory's limits. Processes illuminate ships, rivers, flames, and people, but why are some stable like objects? What sets a stone's slow pace? If you are temporal slices, why should today's reader answer for last year's copied homework? Law, promises, and regret presuppose one you through time. The process picture lights old problems and ignites new ones, as good theories often do.

\section{Today: Theseus's Ship in the Cloud}
The ancient problem wears new clothes nearby. Consider familiar examples.

\textbf{Is your old school still your school?} New heads, campuses, mergers, and names may make alumni deny continuity while principals celebrate an unbroken century. One uses buildings, people, names; the other tradition. This is not mere sentimentality but competing identity criteria. Conservators face the same question in ancient towers mostly repaired through later dynasties. Their refined continuity approach preserves original components and historical information while making additions identifiable and reversible, giving identity a workable craft boundary.

\textbf{A familiar scientific counterclaim is unreliable.} You may hear that every bodily cell renews in seven years, making you materially another person. This apparent support for continuity is too crude. Gut cells change within days, skin within weeks, but many cells, especially most brain neurons, last a lifetime. Your ship replaces some parts often and others never, not all planks every seven years. Correction deepens the question: does identity rest on lasting neurons, relatively stable connections more like structure than material, or Locke's memory? Fact-checking replaces a rumor with better questions.

\textbf{The digital world mass-produces Theseus's ship.} Accounts migrate devices and servers with no original bits remaining, yet remain yours. We instinctively prioritize continuity and structure, not original electrons. Open-source software may retain its name after every original code line changes. AI pushes further: are two copies on different servers one person or two? Is shutting one down death or losing a duplicate? Textbooks cannot settle these new ships, although Plutarch's, Locke's, and Heraclitus's tools remain useful.

\textbf{Consider the classroom itself.} After transfers, changing teachers, and two room moves, graduates still announce Class Three. Neither purely material nor structural, this is more like a speech act: a group declares it remains us. Identity sometimes rests on continued mutual recognition. Nations, teams, families, and bands largely inhabit such recognition. A fourth answer? Perhaps. Metaphysics' list is still growing.

\section{For You: Is the Copy You?}
Return to Geng and the old man, recognizing a bike or insisting on its lost parts. Push their dispute to the limit and answer with reasons.

Suppose future technology scans brain and body atom by atom, making a perfect copy in another city with all your memories, character, and habits. It wakes thinking it has been transported. After scanning, the original you is harmlessly disassembled.

\textbf{Is that copy you?}

Before concluding, lay out the chapter's components again.

Materialists about identity say no: no original atoms traveled; you died and a perfect replica remains. But why does ordinary bodily replacement preserve you, and why measure bicycles and people differently?

Continuity brings trouble too. Memory, personality, and structure continue, apparently preserving you. But suppose the original is \textbf{not} disassembled. Both wake claiming originality. Continuity branches, while you seemingly cannot be two unrelated people.

Conventionalists leave identity to agreed naming. Yet lying inside the scanner as its door closes, would merely a naming question really reassure you?

Process thinking offers a final card. Perhaps you were never a copyable thing, but a process that can continue without being transported, like a melody that cannot be moved to another room, only sounded there anew. Is what sounds the same piece?

There is no standard answer. Athenians arguing over rotten planks 2,400 years ago did not know they rehearsed today's scanner dispute. Metaphysics inspects basic words like thing, same, and exist, usually transparent through usefulness until boundary cases expose cracks.

Next time you say this is still that, pause and ask what warrants the is.

From that moment, you stand at metaphysics' door.
